% GENERATED FILE. Edit canonical lessons, then run npm run generate:books.
% canonical-source-hash: fnv1a64:de5b6524ddb6f42f
% canonical-lessons: HI-C198-sabd, HI-C198-vakya, HI-C198-arth, HI-C198-jnan, HI-C198-vijnan

\chapter{A word, A sentence, A meaning, Knowledge, Science}
\label{ch:hi-a-word-a-sentence-a-meaning-knowledge-science}
\hlchaptermodality{198}

\begin{chapteropening}
\textbf{By the end of this chapter:} \emph{I can say a word, a sentence, a meaning, knowledge and science.}

Complete the last lesson of chapter 198: i can say a word, a sentence, a meaning, knowledge and science.
\end{chapteropening}

\section[\texorpdfstring{\hi{शब्द} (śabd)}{शब्द (śabd)}]{\hi{शब्द} (śabd) — a word}

\label{lesson:HI-C198-sabd}



\begin{quote}
Before the new one: say the Hindi for practice, then the Hindi for a mark.
\end{quote}

\subsection*{You'll want to know: \hi{शब्द}}
\textbf{\hi{शब्द}} — \emph{śabd} — \textquotedblleft{}a word\textquotedblright{}.

\begin{grammarlens}[title={What you've built}]
One more word for this chapter.
\end{grammarlens}

\subsection*{Your turn}
\begin{itemize}
\raggedright
  \item \emph{Say it:} \emph{śabd}
  \item \emph{Say it:} \emph{śabd}, once more
  \item \emph{Say it:} say \emph{śabd}
  \item \emph{Recall:} say the Hindi for practice, then the Hindi for a mark, then say \emph{śabd} again
\end{itemize}

\begin{tcolorbox}[breakable,colback=teal!4,colframe=teal!35!black,title={Before you move on}]
What does \hi{शब्द} mean? (\textquotedblleft{}A word\textquotedblright{}.) Say it once more.
\end{tcolorbox}
\section[\texorpdfstring{\hi{वाक्य} (vākya)}{वाक्य (vākya)}]{\hi{वाक्य} (vākya) — a sentence}

\label{lesson:HI-C198-vakya}



\begin{quote}
Before the new one: say the Hindi for a mark, then the Hindi for a word.
\end{quote}

\subsection*{You'll want to know: \hi{वाक्य}}
\textbf{\hi{वाक्य}} — \emph{vākya} — \textquotedblleft{}a sentence\textquotedblright{}.

\begin{grammarlens}[title={What you've built}]
One more word for this chapter.
\end{grammarlens}

\subsection*{Your turn}
\begin{itemize}
\raggedright
  \item \emph{Say it:} \emph{vākya}
  \item \emph{Say it:} \emph{vākya}, once more
  \item \emph{Say it:} say \emph{vākya}
  \item \emph{Recall:} say the Hindi for a mark, then the Hindi for a word, then say \emph{vākya} again
\end{itemize}

\begin{tcolorbox}[breakable,colback=teal!4,colframe=teal!35!black,title={Before you move on}]
What does \hi{वाक्य} mean? (\textquotedblleft{}A sentence\textquotedblright{}.) Say it once more.
\end{tcolorbox}
\section[\texorpdfstring{\hi{अर्थ} (arth)}{अर्थ (arth)}]{\hi{अर्थ} (arth) — a meaning}

\label{lesson:HI-C198-arth}



\begin{quote}
Before the new one: say the Hindi for a word, then the Hindi for a sentence.
\end{quote}

\subsection*{You'll want to know: \hi{अर्थ}}
\textbf{\hi{अर्थ}} — \emph{arth} — \textquotedblleft{}a meaning\textquotedblright{}.

\begin{grammarlens}[title={What you've built}]
One more word for this chapter.
\end{grammarlens}

\subsection*{Your turn}
\begin{itemize}
\raggedright
  \item \emph{Say it:} \emph{arth}
  \item \emph{Say it:} \emph{arth}, once more
  \item \emph{Say it:} say \emph{arth}
  \item \emph{Recall:} say the Hindi for a word, then the Hindi for a sentence, then say \emph{arth} again
\end{itemize}

\begin{tcolorbox}[breakable,colback=teal!4,colframe=teal!35!black,title={Before you move on}]
What does \hi{अर्थ} mean? (\textquotedblleft{}A meaning\textquotedblright{}.) Say it once more.
\end{tcolorbox}
\section[\texorpdfstring{\hi{ज्ञान} (jñān)}{ज्ञान (jñān)}]{\hi{ज्ञान} (jñān) — knowledge}

\label{lesson:HI-C198-jnan}



\begin{quote}
Before the new one: say the Hindi for a sentence, then the Hindi for a meaning.
\end{quote}

\subsection*{You'll want to know: \hi{ज्ञान}}
\textbf{\hi{ज्ञान}} — \emph{jñān} — \textquotedblleft{}knowledge\textquotedblright{}.

\begin{grammarlens}[title={What you've built}]
One more word for this chapter.
\end{grammarlens}

\subsection*{Your turn}
\begin{itemize}
\raggedright
  \item \emph{Say it:} \emph{jñān}
  \item \emph{Say it:} \emph{jñān}, once more
  \item \emph{Say it:} say \emph{jñān}
  \item \emph{Recall:} say the Hindi for a sentence, then the Hindi for a meaning, then say \emph{jñān} again
\end{itemize}

\begin{tcolorbox}[breakable,colback=teal!4,colframe=teal!35!black,title={Before you move on}]
What does \hi{ज्ञान} mean? (\textquotedblleft{}Knowledge\textquotedblright{}.) Say it once more.
\end{tcolorbox}
\section[\texorpdfstring{\hi{विज्ञान} (vijñān)}{विज्ञान (vijñān)}]{\hi{विज्ञान} (vijñān) — science}

\label{lesson:HI-C198-vijnan}



\begin{quote}
Before the new one: say the Hindi for a meaning, then the Hindi for knowledge.
\end{quote}

\subsection*{You'll want to know: \hi{विज्ञान}}
\textbf{\hi{विज्ञान}} — \emph{vijñān} — \textquotedblleft{}science\textquotedblright{}.

\begin{grammarlens}[title={What you've built}]
One more word for this chapter.
\end{grammarlens}

\subsection*{Your turn}
\begin{itemize}
\raggedright
  \item \emph{Say it:} \emph{vijñān}
  \item \emph{Say it:} \emph{vijñān}, once more
  \item \emph{Say it:} say \emph{vijñān}
  \item \emph{Recall:} say the Hindi for a meaning, then the Hindi for knowledge, then say \emph{vijñān} again
\end{itemize}

\begin{tcolorbox}[breakable,colback=teal!4,colframe=teal!35!black,title={Before you move on}]
What does \hi{विज्ञान} mean? (\textquotedblleft{}Science\textquotedblright{}.) Say it once more.
\end{tcolorbox}
